%% ma: L12-B07-0030
%% lop: 12
%% chuong: 2
%% bai: 7
%% dang: 12.7.7
%% loai: TN4
%% muc: TH
%% dap_an: "A"
%% nguon: "CĐ2.6.pdf giáo viên gửi 27/10/2026 (M07, Lê Minh Khôi – Nguyễn Anh Khôi, Chương II Chủ đề 6 Ứng dụng vào các bài toán thực tế), Câu 10"
%% kien_thuc: "Bài 7 (tọa độ theo các hướng Đông, Tây, Nam, Bắc), Bài 8 (khoảng cách hai điểm)"
%% trang_thai: chinh-thuc
\begin{ex}
Hai chiếc máy bay không người lái cùng bay lên tại một địa điểm. Sau một thời gian bay, chiếc máy bay thứ nhất cách điểm xuất phát về phía Bắc $20$ km và về phía Tây $10$ km, đồng thời cách mặt đất $700$ m. Chiếc máy bay thứ hai cách điểm xuất phát về phía Đông $30$ km và về phía Nam $25$ km, đồng thời cách mặt đất $1$ km. Khoảng cách giữa hai chiếc máy bay lúc này gần nhất với giá trị nào sau đây?
\choice{\True $60{,}21$ km}{$58{,}31$ km}{$65{,}00$ km}{$55{,}40$ km}
\loigiai{Chọn $Ox$ hướng Nam, $Oy$ hướng Đông, $Oz$ hướng lên (đơn vị km): $M_1(-20;-10;0{,}7)$, $M_2(25;30;1)$.

$M_1M_2=\sqrt{45^2+40^2+0{,}3^2}\approx60{,}21$ km.}
\end{ex}
